\documentclass[11pt,a4paper]{report}

% --- Pacchetti Fondamentali ---
\usepackage[utf8]{inputenc}
\usepackage[T1]{fontenc}
\usepackage{lmodern}
\usepackage[italian]{babel}
\usepackage{etoolbox}
\usepackage[margin=1in]{geometry}
\usepackage{hyperref}
\usepackage{xcolor}
\usepackage{listings}
\usepackage{enumitem}
\usepackage{titlesec}
\usepackage{fancyhdr}
\usepackage{cleveref}

% --- Configurazione Header e Footer ---
\pagestyle{fancy}
\fancyhf{}
\fancyhead[L]{\small Programmazione e Tecnologie Web}
\fancyhead[R]{\small Esercitazione Guidata - Lezione 2}
\fancyfoot[C]{\thepage}
\renewcommand{\headrulewidth}{0.4pt}

% --- Configurazione Colori ---
\definecolor{primary}{RGB}{38, 66, 139}
\definecolor{codebg}{RGB}{245, 245, 245}

\hypersetup{
    colorlinks=true,
    linkcolor=primary,
    urlcolor=primary,
}

% --- Configurazione Formattazione Codice ---
\lstset{
    backgroundcolor=\color{codebg},
    basicstyle=\ttfamily\small,
    breaklines=true,
    frame=single,
    rulecolor=\color{lightgray},
    captionpos=b,
    xleftmargin=1em,
    xrightmargin=1em
}

% --- Formattazione Titoli ---
\titleformat{\chapter}[display]
  {\normalfont\bfseries\color{primary}}{\huge Esercitazione \thechapter}{1ex}{\Huge}
\titlespacing*{\chapter}{0pt}{-20pt}{40pt}

\begin{document}

% --- Frontespizio ---
\begin{titlepage}
    \centering
    \vspace*{4cm}
    {\Huge \textbf{\color{primary} Programmazione e Tecnologie Web}} \\[1.5cm]
    {\Large \textbf{Esercitazione Guidata - Lezione 2}} \\[0.5cm]
    {\large Docente: Giuseppe Capasso} \\[1.5cm]
    {\normalsize \textbf{Pubblicazione Web, SSG (Hugo) e SSR (Django)}} \\
    \vfill
    {\normalsize Anno Accademico 2025/2026}
\end{titlepage}

\newpage

% ==========================================
% INTRODUZIONE
% ==========================================
\chapter*{Obiettivi dell'Esercitazione}
In questa sessione pratica impareremo a:
\begin{enumerate}
    \item Configurare un ambiente di lavoro con \textbf{GitHub}.
    \item Creare e pubblicare un sito statico professionale con \textbf{Hugo}.
    \item Sviluppare una semplice applicazione dinamica con \textbf{Django} e \textbf{TailwindCSS}.
    \item Utilizzare Git per gestire due progetti distinti.
\end{enumerate}

% ==========================================
% PARTE 1
% ==========================================
\chapter{Setup Ambiente e GitHub}

\section{Creazione Account GitHub}
Se non ne possiedi gia' uno, vai su \url{https://github.com} e registrati. GitHub sara' il tuo "portfolio" online dove memorizzerai tutto il codice che scriverai durante il corso.

\section{Configurazione Git}
Apri il terminale e assicurati che Git sia configurato correttamente con i tuoi dati:
\begin{lstlisting}[language=bash]
git config --global user.name "Tuo Nome"
git config --global user.email "tua@email.it"
\end{lstlisting}

% ==========================================
% PARTE 2
% ==========================================
\chapter{Web Statico con Hugo}

In questa parte creeremo un sito statico e lo pubblicheremo online.

\section{Installazione di Hugo}
Scarica Hugo per il tuo sistema operativo:
\begin{itemize}
    \item \textbf{Windows}: Scarica l'eseguibile da \url{https://gohugo.io/installation/windows/}.
    \item \textbf{Mac/Linux}: Usa il gestore pacchetti (\texttt{brew install hugo} o \texttt{sudo apt install hugo}).
\end{itemize}

\section{Creazione del Progetto}
Esegui i seguenti comandi nel terminale:
\begin{lstlisting}[language=bash]
# 1. Crea il sito
hugo new site mio-sito-hugo
cd mio-sito-hugo

# 2. Scegli un tema (es. Ananke)
git init
git submodule add https://github.com/theNewDynamic/gohugo-theme-ananke.git themes/ananke
echo "theme = 'ananke'" >> hugo.toml

# 3. Crea il primo post
hugo new posts/benvenuto.md
\end{lstlisting}

\section{Pubblicazione}
Crea un nuovo repository su GitHub chiamato \textbf{mio-sito-hugo}, poi collega la tua cartella locale e fai il push:
\begin{lstlisting}[language=bash]
git add .
git commit -m "Primo sito Hugo"
git branch -M main
git remote add origin https://github.com/TUO_UTENTE/mio-sito-hugo.git
git push -u origin main
\end{lstlisting}

\textbf{Task}: Vai su Cloudflare Pages e collega questo repository per vedere il sito online!

% ==========================================
% PARTE 3
% ==========================================
\chapter{Web Dinamico con Django}

Passiamo ora al Server-Side Rendering (SSR). Prima di installare Django, e' fondamentale creare un \textbf{Ambiente Virtuale (Virtual Env)}. Questo serve a isolare le librerie del progetto ed evitare conflitti tra diverse versioni di Python.

\section{Creazione del Virtual Environment}
Crea una cartella per il tuo progetto e, al suo interno, genera l'ambiente virtuale:

\subsection{Per utenti Windows}
Apri il terminale (PowerShell o Prompt dei comandi) e scrivi:
\begin{lstlisting}[language=bash]
# 1. Crea l'ambiente
python -m venv venv

# 2. Attivalo
.\venv\Scripts\activate
\end{lstlisting}

\subsection{Per utenti macOS e Linux}
Apri il terminale e scrivi:
\begin{lstlisting}[language=bash]
# 1. Crea l'ambiente
python3 -m venv venv

# 2. Attivalo
source venv/bin/activate
\end{lstlisting}

\textbf{Nota}: Una volta attivato, vedrai la scritta \texttt{(venv)} all'inizio della riga del tuo terminale. Da questo momento in poi, ogni comando \texttt{pip} installera' i pacchetti solo per questo progetto.

\section{Installazione Django}
Con l'ambiente virtuale \textbf{attivato}, installa Django:
\begin{lstlisting}[language=bash]
pip install django
\end{lstlisting}

\section{Creazione Applicazione}
Creeremo un progetto Django e una singola vista che inietta dati in un template TailwindCSS.
\begin{lstlisting}[language=bash]
# Crea il progetto
django-admin startproject mio_progetto_django
cd mio_progetto_django

# Avvia il server locale
python manage.py runserver
\end{lstlisting}

\section{Configurazione Template SSR}
Crea una cartella \texttt{templates/} e inserisci il file \texttt{django\_tailwind.html} fornito dal docente. 

\textbf{Esercizio}: Modifica il file \texttt{views.py} del tuo progetto (o creane uno) per passare un dizionario di dati (context) al template, includendo il tuo nome e una lista di "cose da fare".

\begin{lstlisting}[language=Python]
# Esempio di logica nella view
from django.shortcuts import render

def home(request):
    context = {
        'user_name': 'Giuseppe',
        'tasks': [
            {'title': 'Imparare Django', 'done': False, 'priority': 'high'},
            {'title': 'Fare la spesa', 'done': True, 'priority': 'low'},
        ]
    }
    return render(request, 'django_tailwind.html', context)
\end{lstlisting}

\section{Secondo Repository}
Crea un secondo repository su GitHub chiamato \textbf{mio-progetto-django} e carica tutto il codice. Ricorda di creare un file \texttt{.gitignore} per non caricare la cartella \texttt{\_\_pycache\_\_}!

\end{document}
